% GENERATED FILE. Edit canonical lessons, then run npm run generate:books.
% canonical-source-hash: fnv1a64:54e40cf18eb05d25
% canonical-lessons: AR-C342-turath, AR-C342-makhtuta, AR-C342-naqsh, AR-C342-laqab, AR-C342-ikhtaraa

\chapter{Heritage and inventions}
\label{ch:ar-heritage-and-inventions}
\hlchaptermodality{342}

\begin{chapteropening}
\textbf{By the end of this chapter:} \emph{I can say heritage, a manuscript, an inscription, a title and to invent.}

Complete the last lesson of chapter 342: I can say heritage, a manuscript, an inscription, a title and to invent.
\end{chapteropening}

\section[\texorpdfstring{\ar{تراث} (turāth)}{تراث (turāth)}]{\ar{تراث} (turāth) — heritage}

\label{lesson:AR-C342-turath}



\begin{quote}
Before the new one: say the Arabic for a treaty, then the Arabic for to betray.
\end{quote}

\subsection*{You'll want to know: \ar{تراث}}
\textbf{\ar{تراث}} — \emph{turāth} — \textquotedblleft{}heritage\textquotedblright{}.

\begin{quote}
\ar{تحافظ} \ar{القرية} \ar{على} \ar{تراثها} \ar{القديم}.
\end{quote}

\begin{quote}
\emph{tuḥāfiẓu al-qarya ʿalā turāthihā al-qadīm.}
\end{quote}

\begin{quote}
The village takes care of its old heritage.
\end{quote}

\begin{cousinweb}
\textbf{\ar{تراث}} (\emph{turāth}) is built on the native Arabic root \emph{w-r-th} (Lane), \textquotedblleft{}to inherit\textquotedblright{}. Lane explains \emph{turāth} as \emph{wurāth} with \emph{t} in place of the \emph{w}: what a people inherits.
\end{cousinweb}

\begin{grammarlens}[title={What you've built}]
One more word for telling what happened.
\end{grammarlens}

\subsection*{Your turn}
\begin{itemize}
\raggedright
  \item \emph{Say it:} \emph{turāth}
  \item \emph{Say it:} \emph{turāth}, once more
  \item \emph{Say it:} \emph{tuḥāfiẓu al-qarya ʿalā turāthihā al-qadīm.} — The village takes care of its old heritage.
  \item \emph{Recall:} say the Arabic for a treaty, then the Arabic for to betray, then say \emph{turāth} again
\end{itemize}

\begin{tcolorbox}[breakable,colback=teal!4,colframe=teal!35!black,title={Before you move on}]
What does \ar{تراث} mean? (\textquotedblleft{}Heritage\textquotedblright{}.) Say it once more.
\end{tcolorbox}
\section[\texorpdfstring{\ar{مخطوطة} (makhṭūṭa)}{مخطوطة (makhṭūṭa)}]{\ar{مخطوطة} (makhṭūṭa) — a manuscript}

\label{lesson:AR-C342-makhtuta}



\begin{quote}
Before the new one: say the Arabic for to betray, then the Arabic for heritage.
\end{quote}

\subsection*{You'll want to know: \ar{مخطوطة}}
\textbf{\ar{مخطوطة}} — \emph{makhṭūṭa} — \textquotedblleft{}a manuscript\textquotedblright{}.

\begin{quote}
\ar{وجد} \ar{الباحث} \ar{مخطوطة} \ar{قديمة} \ar{في} \ar{المكتبة}.
\end{quote}

\begin{quote}
\emph{wajada al-bāḥith makhṭūṭa qadīma fī al-maktaba.}
\end{quote}

\begin{quote}
The researcher found an old manuscript in the library.
\end{quote}

Plural \ar{مخطوطات}.

\begin{cousinweb}
\textbf{\ar{مخطوطة}} (\emph{makhṭūṭa}) is built on the native Arabic root \emph{kh-ṭ-ṭ} (Lane), \textquotedblleft{}to draw a line; to write\textquotedblright{}, the root of \emph{khaṭṭ}, handwriting. \emph{Makhṭūṭa} is \textquotedblleft{}a thing written by hand\textquotedblright{}.
\end{cousinweb}

\begin{grammarlens}[title={What you've built}]
One more word for telling what happened.
\end{grammarlens}

\subsection*{Your turn}
\begin{itemize}
\raggedright
  \item \emph{Say it:} \emph{makhṭūṭa}
  \item \emph{Say it:} \emph{makhṭūṭa}, once more
  \item \emph{Say it:} \emph{wajada al-bāḥith makhṭūṭa qadīma fī al-maktaba.} — The researcher found an old manuscript in the library.
  \item \emph{Recall:} say the Arabic for to betray, then the Arabic for heritage, then say \emph{makhṭūṭa} again
\end{itemize}

\begin{tcolorbox}[breakable,colback=teal!4,colframe=teal!35!black,title={Before you move on}]
What does \ar{مخطوطة} mean? (\textquotedblleft{}A manuscript\textquotedblright{}.) Say it once more.
\end{tcolorbox}
\section[\texorpdfstring{\ar{نقش} (naqsh)}{نقش (naqsh)}]{\ar{نقش} (naqsh) — an inscription; a carving}

\label{lesson:AR-C342-naqsh}



\begin{quote}
Before the new one: say the Arabic for heritage, then the Arabic for a manuscript.
\end{quote}

\subsection*{You'll want to know: \ar{نقش}}
\textbf{\ar{نقش}} — \emph{naqsh} — \textquotedblleft{}an inscription; a carving\textquotedblright{}.

\begin{quote}
\ar{قرأ} \ar{العالم} \ar{نقشا} \ar{على} \ar{الحجر}.
\end{quote}

\begin{quote}
\emph{qaraʾa al-ʿālim naqshan ʿalā al-ḥajar.}
\end{quote}

\begin{quote}
The scholar read an inscription on the stone.
\end{quote}

Plural: \ar{نقوش}, \emph{nuqūsh}.

\begin{cousinweb}
\textbf{\ar{نقش}} (\emph{naqsh}) is built on the native Arabic root \emph{n-q-sh} (Lane), \textquotedblleft{}to paint or engrave a design\textquotedblright{}; the taught \emph{nāqasha}, to discuss, is picking at a matter point by point.
\end{cousinweb}

\begin{grammarlens}[title={What you've built}]
One more word for telling what happened.
\end{grammarlens}

\subsection*{Your turn}
\begin{itemize}
\raggedright
  \item \emph{Say it:} \emph{naqsh}
  \item \emph{Say it:} \emph{naqsh}, once more
  \item \emph{Say it:} \emph{qaraʾa al-ʿālim naqshan ʿalā al-ḥajar.} — The scholar read an inscription on the stone.
  \item \emph{Recall:} say the Arabic for heritage, then the Arabic for a manuscript, then say \emph{naqsh} again
\end{itemize}

\begin{tcolorbox}[breakable,colback=teal!4,colframe=teal!35!black,title={Before you move on}]
What does \ar{نقش} mean? (\textquotedblleft{}An inscription; a carving\textquotedblright{}.) Say it once more.
\end{tcolorbox}
\section[\texorpdfstring{\ar{لقب} (laqab)}{لقب (laqab)}]{\ar{لقب} (laqab) — a title; a nickname}

\label{lesson:AR-C342-laqab}



\begin{quote}
Before the new one: say the Arabic for a manuscript, then the Arabic for an inscription.
\end{quote}

\subsection*{You'll want to know: \ar{لقب}}
\textbf{\ar{لقب}} — \emph{laqab} — \textquotedblleft{}a title; a nickname\textquotedblright{}.

\begin{quote}
\ar{كان} \ar{لقب} \ar{جدي} \ar{في} \ar{القرية} \ar{الحكيم}.
\end{quote}

\begin{quote}
\emph{kāna laqab jaddī fī al-qarya al-ḥakīm.}
\end{quote}

\begin{quote}
My grandfather’s nickname in the village was the Wise One.
\end{quote}

Plural: \ar{ألقاب}, \emph{alqāb}.

In sport: \ar{لقب} \ar{البطولة}, the title.

\begin{cousinweb}
\textbf{\ar{لقب}} (\emph{laqab}) is built on the native Arabic root \emph{l-q-b} (Lane), where \emph{laqab} is a name given to a person besides his own, for praise or in jest.
\end{cousinweb}

\begin{grammarlens}[title={What you've built}]
One more word for telling what happened.
\end{grammarlens}

\subsection*{Your turn}
\begin{itemize}
\raggedright
  \item \emph{Say it:} \emph{laqab}
  \item \emph{Say it:} \emph{laqab}, once more
  \item \emph{Say it:} \emph{kāna laqab jaddī fī al-qarya al-ḥakīm.} — My grandfather’s nickname in the village was the Wise One.
  \item \emph{Recall:} say the Arabic for a manuscript, then the Arabic for an inscription, then say \emph{laqab} again
\end{itemize}

\begin{tcolorbox}[breakable,colback=teal!4,colframe=teal!35!black,title={Before you move on}]
What does \ar{لقب} mean? (\textquotedblleft{}A title; a nickname\textquotedblright{}.) Say it once more.
\end{tcolorbox}
\section[\texorpdfstring{\ar{اخترع} (ikhtaraʿa)}{اخترع (ikhtaraʿa)}]{\ar{اخترع} (ikhtaraʿa) — to invent}

\label{lesson:AR-C342-ikhtaraa}



\begin{quote}
Before the new one: say the Arabic for an inscription, then the Arabic for a title.
\end{quote}

\subsection*{You'll want to know: \ar{اخترع}}
\textbf{\ar{اخترع}} — \emph{ikhtaraʿa} — \textquotedblleft{}to invent\textquotedblright{}.

\begin{quote}
\ar{اخترع} \ar{العالم} \ar{آلة} \ar{جديدة}.
\end{quote}

\begin{quote}
\emph{ikhtaraʿa al-ʿālim āla jadīda.}
\end{quote}

\begin{quote}
The scientist invented a new machine.
\end{quote}

Present: \ar{يخترع}, \emph{yakhtariʿ}.

An invention is \ar{اختراع}.

\begin{cousinweb}
\textbf{\ar{اخترع}} (\emph{ikhtaraʿa}) is built on the native Arabic root \emph{kh-r-ʿ} (Lane), \textquotedblleft{}to split, to cleave\textquotedblright{}. It is Form VIII, \textquotedblleft{}to bring a thing out new, without a model\textquotedblright{}.
\end{cousinweb}

\begin{grammarlens}[title={What you've built}]
One more word for telling what happened.
\end{grammarlens}

\subsection*{Your turn}
\begin{itemize}
\raggedright
  \item \emph{Say it:} \emph{ikhtaraʿa}
  \item \emph{Say it:} \emph{ikhtaraʿa}, once more
  \item \emph{Say it:} \emph{ikhtaraʿa al-ʿālim āla jadīda.} — The scientist invented a new machine.
  \item \emph{Recall:} say the Arabic for an inscription, then the Arabic for a title, then say \emph{ikhtaraʿa} again
\end{itemize}

\begin{tcolorbox}[breakable,colback=teal!4,colframe=teal!35!black,title={Before you move on}]
What does \ar{اخترع} mean? (\textquotedblleft{}To invent\textquotedblright{}.) Say it once more.
\end{tcolorbox}
